\section{Introduction}
% For modern-day DNN architectures, billions to trillions of parameters of large models have made it increasingly challenging for developers. For example, Large Language Models (LLM) like BERT or GPT-3 have billions of parameters and require more than 100 Gigabytes of GPU memory to store activation during one training epoch. The memory requirements of these models prevent ML developers from training them on memory-limited devices. Without the help of advanced scaling techniques such as ZeRO, Rematerialization, and hybrid parallelism, ordinary users cannot determine whether a customized model could fit into GPU memory during execution, given certain information of \{\textit{BatchSize, SeqLen, ModelConfigs}\}. Even if the device has adequate memory to serve the large model, it still requires non-trivial time to profile the runtime information. This will confine the deep learning system to generate static optimization plans only, or requires researchers to devise dynamic strategies that are often heuristic and sub-optimal.

% Therefore, A feasible solution is to trace the model's computation graph with meta-information without executing the actual kernels on real devices. Thanks to the new feature of PyTorch 1.12 release, symbolic tracing of \texttt{torch.fx} made it possible to capture the program of a model without real execution. Another exciting feature of \texttt{\_\_torch\_dispatch\_\_} enabled user to allocate and propagate a \texttt{torch.Tensor} object on \texttt{device='meta'}. This is highly efficient because the execution works the same as tensor shape propagation. Moreover, since shape propagation only requires a tensor's shape and dtype information, it is not necessary to allocate any physical memory on virtual devices. In other words, the total execution of profiling done in seconds, and there are no requirements for storing actual data on GPU memory.

% In summary, the main contributions of this paper are:
% \begin{itemize}
%     \item An implementation of an easy-to-use meta-profiler to generate memory and
% FLOPs statistics of an arbitrary PyTorch model.
%     \item An in-depth discussion of memory allocation mechanism for \texttt{torch.autograd}
% \end{itemize}

% Foundation models~\cite{bommasani2021opportunities} have emerged as a cutting-edge paradigm in the field of AI. Neural networks like BERT~\cite{bert}, GPT-3~\cite{gpt-3} and Vision Transformer~\cite{https://doi.org/10.48550/arxiv.2010.11929} contain large number of parameters and trained on huge datasets. They are capable to learn potential high-dimensional features and show better performance as the number of parameters grows. Together with the blooming of powerful hardware, the AI community has seen a trend of models getting larger. As this trend goes on, it seems that single GPU could not catch up with such wild demand of fast memory. Without any optimization, training a model of 10 billion parameters can cost more than 80 GB memory even with one sample per batch. Therefore, large models in modern days are trained in distributed manner, systems like GShard~\cite{gshard}, FairScale~\cite{FairScale2021}, Megatron-LM~\cite{shoeybi2019megatron} and DeepSpeed~\cite{deepspeed} provide strategies on distributed training to train larger model on limited devices. 

% % Research has demonstrated that as the number of parameters increases, the performance of foundation models improves. 
% % However, owing to the vast number of parameters and computation involved, these models must be trained in a distributed manner, and usually requires activation checkpoint to alleviate stress on GPU memory.
% % This process typically necessitates a significant amount of specialized engineering work, usually adapting the model to the cluster configuration in a distributed manner~\cite{bian2021colossal}.
% % This can be a complex and laborious task, requiring a thorough understanding of both the model and the cluster environment.

% % Recent advancements in machine learning systems have led to the emergence of automatic parallel systems, which aim to simplify the process of distributed training for foundational models. 
% % These systems utilize an automated methodology for determining the most efficient parallel execution plan based on the underlying hardware configurations. The automatic parallel system distributes the model across multiple devices by utilizing two orthogonal parallel strategies: intra-op parallelism and inter-op parallelism. The first one is more complex and encompasses a larger search space. In view of this, this paper aims to specifically focus on the design of intra-op parallelism.


% However, it still requires heavy engineering work on customized models to find the optimal execution plan with SPMD style intra-op parallelism~\cite{xu2021gspmd}. 
% First, as each dimension of the operator tensor could be partitioned and devices can be viewed as an arbitrary N-D mesh, the combination of two factors leads to a large search space.
% Second, the output of an operator is the input of another one, which may require a completely different sharding spec.
% Since it is almost impossible to exhaustively implement the transformation between two sharding specs, 
% the existing method requires that the number of dimensions of the device mesh and tensor partition is less than 2.
% Third, the search for the best intra-op parallelism can be interfered with by other optimizations.
% For example, activation recomputation~\cite{act-ckpt} and model data offloading strategies~\cite{ren2021zero, fang2022parallel} can reduce memory consumption.
% Whether and how these strategies are used can lead to a completely different search space for intra-op parallelism. Therefore, building an automatic parallel systems is a significant task to simplify the process to find the optimal execution plan.

% Current automatic intra-op parallelism solutions, such as Alpa~\cite{zheng2022alpa}, Flexflow~\cite{jia2019beyond} still have some limitations. 
% First, activation recomputation has already become a must-have recipe for foundation model training, but the existing methods do not take it into consideration.
% second, the generalization for new hardware configurations is insufficient. Existing work can only organize the devices into a 2D mesh, and cannot work for more complex typologies.

% We have implemented an automatic intra-op parallel system based on PyTorch~\cite{paszke2019pytorch}, which compiles the serial model training code into efficient parallel execution code. 
% Different from the other works, our system is able to automatically generate the most efficient activation recomputation scheme for an execution plan to meet the device memory budget.
% To achieve this goal, it collects information from both hardware and DNN computation graph in a profiling phase, 
% adopts a hierarchical optimization method to find out the optimal execution plan, 
% and generates the runtime execution code. Specifically, our contributions include:

% \begin{itemize}
%     \item We jointly optimize the intra-op parallelism and activation recomputation methods ahead of time. 
%     \item A heuristic algorithm to convert the tensor layout of different sharding specs is developed to support tensor layout converting in an arbitrary dimension device mesh.
%     \item We implement a profiler which could collect both computing and memory overhead in the fine-grained computation graph using meta-execution.
% \end{itemize}


Foundation models~\cite{bommasani2021opportunities} have emerged as a leading paradigm in the field of AI, with neural networks such as BERT~\cite{bert}, GPT-3~\cite{gpt-3}, and Vision Transformer~\cite{https://doi.org/10.48550/arxiv.2010.11929} containing a large number of parameters and trained on massive datasets. As the number of parameters grows, these models show improved performance, and with the advent of powerful hardware, there has been a trend toward building increasingly larger models. However, this trend has resulted in a situation where a single GPU is unable to meet the high demand for fast memory. Without any optimization, training a model of 10 billion parameters can cost more than 80 GB of memory, even with just one sample per batch. As a result, large models in modern times are trained in a distributed manner, with systems such as GShard~\cite{gshard}, FairScale~\cite{FairScale2021}, Megatron-LM~\cite{shoeybi2019megatron}, and DeepSpeed~\cite{deepspeed} providing strategies for distributed training on limited devices.

However, distributed training of large models requires significant engineering effort to find the optimal execution plan with Single Program Multiple Data (SPMD) style intra-op parallelism~\cite{xu2021gspmd}. This is due to the large search space resulting from the combination of two factors: each dimension of the operator tensor can be partitioned, and devices can be viewed as an arbitrary N-dimensional mesh. Furthermore, the output of one operator is the input of another, which may require a completely different sharding specification. Since it is almost impossible to implement the transformation between two sharding specifications exhaustively, existing methods require that the number of dimensions of the device mesh and tensor partition is less than or equal to 2. Additionally, the search for the best intra-op parallelism can be interfered with by other optimizations, such as activation recomputation~\cite{act-ckpt} and model data offloading strategies~\cite{ren2021zero, fang2022parallel}, which can reduce memory consumption and lead to a completely different search space for intra-op parallelism. Therefore, developing an automatic parallel system is a significant but challenging task to simplify the process of finding the optimal execution plan.

Current automatic intra-op parallelism solutions, such as Alpa~\cite{zheng2022alpa} and Flexflow~\cite{jia2019beyond}, still have some limitations. For example, activation recomputation has become a must-have recipe for foundation model training, but existing methods do not take it into consideration. Additionally, the generalization for new hardware configurations is insufficient, as existing work can only organize devices into a 2D mesh and cannot work for more complex typologies.

In this paper, we present an automatic intra-op parallel system based on PyTorch~\cite{paszke2019pytorch}, which compiles serial model code into efficient parallel execution code. Besides parallelization, our system can automatically generate the most efficient activation recomputation scheme for an execution plan to meet the device memory budget. To achieve this, our system collects information from both hardware and  computation graphs in a profiling phase, adopts a hierarchical optimization method to find the optimal execution plan, and generates runtime execution code. Specifically, our contributions include:

\begin{itemize}
    \item We jointly optimize the intra-op parallelism and activation recomputation methods ahead of time. 
    \item A heuristic algorithm to convert the tensor layout of different sharding specs is developed to support efficient tensor layout conversion in an arbitrary dimension device mesh.
    \item We implement a profiler that could collect both computing and memory overhead in the fine-grained computation graph using meta-execution.
\end{itemize}